% Cible : ARIEL_LOGMINER_MEMOIRE_FINAL/chapters/annexes.tex
% Passage à remplacer : de « \section{Matrice résultats--preuves} » jusqu'avant « \begin{longtable} »
% Agent : 6 — Reproductibilité, artefacts et déclaration IA
% Sources : experiments/phase_revision_20261104/{README.md,MANIFEST_SHA256.csv} ;
%           chapitre 7 proposé ; RETOURS_LECTEURS.md, lecteur 1 E.2.

\section{Matrice résultats--preuves}

La matrice ci-dessous fournit un point d'entrée vers les artefacts qui étayent les
résultats du mémoire. Les chemins complets sont relatifs à la racine du dépôt. Dans le
paquet \texttt{experiments/phase\_revision\_20261104}, ils sont abrégés en commençant au
sous-répertoire de campagne, par exemple \texttt{agent\_2\_hdfs/}. Le fichier
\texttt{MANIFEST\_SHA256.csv} de ce paquet permet d'en contrôler le chemin, la taille et
l'empreinte SHA-256.

Cette matrice établit la traçabilité des valeurs, non leur reproduction indépendante.
Une reproduction exige également le script, l'entrée, les dépendances, les paramètres et
les graines associés. Les jeux publics qui ne sont pas redistribués doivent être obtenus
depuis leur source documentée et vérifiés avant l'exécution.
